\documentclass[../../main.tex]{subfiles}
\begin{document}

\section{曲线的速度向量}\label{section:ch03-曲线的速度向量}

$\mathbb{R}^{n}$ 中一条光滑参数化曲线的速度，在初等微积分中是熟悉的。它正是那个分量为该曲线各分量函数导数的向量。本节中我们把这个概念推广到流形中的曲线。

若 $M$ 是一个带边或无边流形，我们定义 $M$ 中的一条曲线为一个连续映射 $\gamma:J\to M$，其中 $J\subset\mathbb{R}$ 是一个区间。（大多数时候我们将感兴趣于定义域为开区间的曲线，但为某些目的，允许 $J$ 有一个或两个端点是有用的；在那种情形下定义经微小修改后全部有意义，或者通过把 $J$ 看作一个带边流形，或者通过把导数解释为单侧导数。）注意本书中术语曲线总是指从一个区间到 $M$ 的映射（一条参数化曲线），而不仅仅是 $M$ 中一个点集。

现在设 $M$ 是一个光滑流形，仍然带边或无边。我们对切空间的定义自然地导向速度向量的一种解释：给定一条光滑曲线 $\gamma:J\to M$ 和 $t_{0}\in J$，我们定义 $\gamma$ 在 $t_{0}$ 处的速度（图~\cref{figure:ch03-3.9}），记作 $\gamma'(t_{0})$，为向量
\begin{align*}
\gamma'(t_{0})=\mathrm{d}\gamma\left(\left.\frac{\mathrm{d}}{\mathrm{d}t}\right|_{t_{0}}\right)\in T_{\gamma(t_{0})}M,
\end{align*}
其中 $\mathrm{d}/\mathrm{d}t|_{t_{0}}$ 是 $T_{t_{0}}\mathbb{R}$ 中的标准坐标基向量。（正如在普通微积分中那样，当流形是 $1$ 维时习惯用 $\mathrm{d}/\mathrm{d}t$ 而不是 $\partial/\partial t$。）速度的其他常见记号是
\begin{align*}
\dot{\gamma}(t_{0}),\qquad \frac{\mathrm{d}\gamma}{\mathrm{d}t}(t_{0}),\qquad \frac{\mathrm{d}\gamma}{\mathrm{d}t}\bigg|_{t=t_{0}}.
\end{align*}

\begin{figure}[ht]
\centering
\includegraphics[width=0.7\textwidth]{fig-03-9.png}
\caption{曲线的速度}\label{figure:ch03-3.9}
\end{figure}

这个切向量对函数的作用为
\begin{align*}
\gamma'(t_{0})f=\mathrm{d}\gamma\left(\left.\frac{\mathrm{d}}{\mathrm{d}t}\right|_{t_{0}}\right)f
=\left.\frac{\mathrm{d}}{\mathrm{d}t}\right|_{t_{0}}(f\circ\gamma)=(f\circ\gamma)'(t_{0}).
\end{align*}
换言之，$\gamma'(t_{0})$ 是在 $\gamma(t_{0})$ 处、通过沿 $\gamma$ 取一个函数的导数而得到的导子。（若 $t_{0}$ 是 $J$ 的一个端点，这仍然成立，只要我们把关于 $t$ 的导数解释为单侧导数，或等价地解释为 $f\circ\gamma$ 到 $\mathbb{R}$ 一个开子集的任意光滑延拓的导数。）

现在设 $(U,\varphi)$ 是一个坐标函数为 $(x^{i})$ 的光滑坐标图。若 $\gamma(t_{0})\in U$，则至少对 $t$ 充分接近 $t_{0}$，我们可以把 $\gamma$ 的坐标表示写成 $\gamma(t)=(\gamma^{1}(t),\dots,\gamma^{n}(t))$，于是微分的坐标公式给出
\begin{align*}
\gamma'(t_{0})=\frac{\mathrm{d}\gamma^{i}}{\mathrm{d}t}(t_{0})\left.\frac{\partial}{\partial x^{i}}\right|_{\gamma(t_{0})}.
\end{align*}
这意味着 $\gamma'(t_{0})$ 本质上由与欧几里得空间中相同的公式给出：它是那个在坐标基中分量为 $\gamma$ 各分量函数导数的切向量。

下一个命题表明，流形上的每个切向量都是某条曲线的速度向量。这给出了另一种、且多少更几何的思考切丛的方式：它正是 $M$ 中光滑曲线的所有速度向量所成的集合。

\begin{proposition}\label{proposition:ch03-切向量是速度}
设 $M$ 是一个带边或无边光滑流形，$p\in M$。每个 $v\in T_{p}M$ 都是 $M$ 中某条光滑曲线的速度。
\end{proposition}

\begin{proof}
首先设 $p\in\mathrm{Int}\,M$（这包括 $\partial M=\varnothing$ 的情形）。令 $(U,\varphi)$ 为以 $p$ 为中心的一个光滑坐标图，并用坐标基把 $v$ 写成 $v=v^{i}\partial/\partial x^{i}|_{p}$。对充分小的 $\varepsilon>0$，令 $\gamma:(-\varepsilon,\varepsilon)\to U$ 为坐标表示为
\begin{align}\label{align:ch03-径向曲线}
\gamma(t)=(tv^{1},\dots,tv^{n})
\end{align}
的曲线。（记住，这实际上意味着 $\gamma(t)=\varphi^{-1}(tv^{1},\dots,tv^{n})$。）这是一条以 $\gamma(0)=p$ 为始点的光滑曲线，上面的计算表明 $\gamma'(0)=v^{i}\partial/\partial x^{i}|_{\gamma(0)}=v$。

现在设 $p\in\partial M$。令 $(U,\varphi)$ 为以 $p$ 为中心的一个光滑边界坐标图，并像前面那样把 $v$ 写成 $v=v^{i}\partial/\partial x^{i}|_{p}$。我们希望让 $\gamma$ 为坐标表示是~\eqref{align:ch03-径向曲线} 的曲线，但这个公式仅当 $tv^{n}\geqslant 0$ 时才代表 $M$ 的一个点。我们可以通过适当限制 $\gamma$ 的定义域来适应这一要求：若 $v^{n}=0$，我们像前面那样定义 $\gamma:(-\varepsilon,\varepsilon)\to U$；若 $v^{n}>0$，我们令定义域为 $[0,\varepsilon)$；若 $v^{n}<0$，我们令定义域为 $(-\varepsilon,0]$。在每种情形下，$\gamma$ 都是 $M$ 中一条以 $\gamma(0)=p$、$\gamma'(0)=v$ 的光滑曲线。
\end{proof}

下一个命题表明，速度向量在与光滑映射复合时表现良好。

\begin{proposition}[复合曲线的速度]\label{proposition:ch03-复合曲线速度}
设 $F:M\to N$ 是一个光滑映射，$\gamma:J\to M$ 是一条光滑曲线。对任意 $t_{0}\in J$，复合曲线 $F\circ\gamma:J\to N$ 在 $t=t_{0}$ 处的速度由
\begin{align*}
(F\circ\gamma)'(t_{0})=\mathrm{d}F(\gamma'(t_{0}))
\end{align*}
给出。
\end{proposition}

\begin{proof}
只需回到曲线速度的定义：
\begin{align*}
(F\circ\gamma)'(t_{0})=\mathrm{d}(F\circ\gamma)\left(\left.\frac{\mathrm{d}}{\mathrm{d}t}\right|_{t_{0}}\right)
=\mathrm{d}F\circ\mathrm{d}\gamma\left(\left.\frac{\mathrm{d}}{\mathrm{d}t}\right|_{t_{0}}\right)
=\mathrm{d}F(\gamma'(t_{0})).
\end{align*}
\end{proof}

表面上看，前一命题告诉我们如何用微分计算复合曲线的速度。然而，反过来使用它常常更有用，把它作为计算微分的一种精简方式。设 $F:M\to N$ 是一个光滑映射，我们需要计算某点 $p\in M$ 处的微分 $\mathrm{d}F_{p}$。我们可以对任意 $v\in T_{p}M$ 计算 $\mathrm{d}F_{p}(v)$：选一条初始切向量为 $v$ 的光滑曲线 $\gamma$，然后把命题~\cref{proposition:ch03-复合曲线速度} 应用于复合曲线 $F\circ\gamma$。下一个推论总结了结果。

\begin{corollary}[用速度向量计算微分]\label{corollary:ch03-用速度算微分}
设 $F:M\to N$ 是一个光滑映射，$p\in M$，$v\in T_{p}M$。则
\begin{align*}
\mathrm{d}F_{p}(v)=(F\circ\gamma)'(0)
\end{align*}
对任意光滑曲线 $\gamma:J\to M$（满足 $0\in J$、$\gamma(0)=p$、$\gamma'(0)=v$）成立。
\end{corollary}

这个推论常常给出 $\mathrm{d}F$ 的一种简洁得多的计算，特别是当 $F$ 以某种非显式坐标表示的形式给出时。我们在后面各章将看到这种技巧的许多例子。

\end{document}